% =========================================================================
% SECCIONES PRELIMINARES: ÍNDICES, ABREVIATURAS, SÍMBOLOS Y GLOSARIO
% =========================================================================

% 1. ÍNDICE GENERAL
\newpage
\phantomsection
\addcontentsline{toc}{chapter}{Índice general}
\tableofcontents 

% 2. ÍNDICE DE FIGURAS
\newpage
\phantomsection
\addcontentsline{toc}{chapter}{Índice de figuras}
\renewcommand{\listfigurename}{Índice de figuras}
\listoffigures

% 3. ÍNDICE DE TABLAS
\newpage
\phantomsection
\addcontentsline{toc}{chapter}{Índice de tablas}
\renewcommand{\listtablename}{Índice de tablas}
\renewcommand{\tablename}{Tabla}
\listoftables

% 4. ÍNDICE DE ALGORITMOS
\newpage
\phantomsection
\addcontentsline{toc}{chapter}{Índice de algoritmos}
\listof{algoritmo}{Índice de algoritmos}

% 5. LISTA DE ABREVIATURAS Y ACRÓNIMOS
\newpage
\phantomsection
\addcontentsline{toc}{chapter}{Lista de abreviaturas}
\chapter*{Lista de abreviaturas}

\begin{table}[H]
\centering
\small
\begin{tabularx}{\textwidth}{@{}p{3.2cm}X@{}}
\toprule
\rowcolor{tcAzul!12}
\textbf{Abreviatura / Sigla} & \textbf{Significado y Descripción Técnica} \\
\midrule
\textbf{AMLQ} & Academia Mayor de la Lengua Quechua (Cusco, Perú). \\
\textbf{BLEU} & \textit{Bilingual Evaluation Understudy} (métrica de concordancia de n-gramas léxicos). \\
\textbf{chrF++} & \textit{Character $n$-gram F-score with word $n$-grams}, métrica de similitud léxica a nivel de caracteres y palabras. \\
\textbf{CPU} & \textit{Central Processing Unit} (unidad central de procesamiento). \\
\textbf{DSR} & \textit{Design Science Research} (marco metodológico de investigación en ciencias del diseño). \\
\textbf{FAISS} & \textit{Facebook AI Similarity Search} (biblioteca de búsqueda vectorial densa eficiente). \\
\textbf{GPU} & \textit{Graphics Processing Unit} (unidad de procesamiento gráfico acelerado). \\
\textbf{HE} & Hipótesis Específica (HE1, HE2, HE3). \\
\textbf{HG} & Hipótesis General de la investigación. \\
\textbf{LLM} & \textit{Large Language Model} (modelo de lenguaje de gran escala). \\
\textbf{LoRA} & \textit{Low-Rank Adaptation} (técnica de adaptación mediante descomposición de bajo rango). \\
\textbf{ODS} & Objetivos de Desarrollo Sostenible de las Naciones Unidas (ODS 9 y ODS 10). \\
\textbf{OE} & Objetivo Específico (OE1, OE2, OE3). \\
\textbf{OG} & Objetivo General de la investigación. \\
\textbf{P@k} & \textit{Precision at k}: proporción de fragmentos pertinentes entre los primeros $k$ recuperados. \\
\textbf{PE} & Problema Específico (PE1, PE2, PE3). \\
\textbf{PG} & Problema General de la investigación. \\
\textbf{PLN} & Procesamiento de Lenguaje Natural (equivalente a NLP en inglés). \\
\textbf{QLoRA} & \textit{Quantized Low-Rank Adaptation} (adaptadores de bajo rango sobre pesos base en 4 bits). \\
\textbf{RAG} & \textit{Retrieval-Augmented Generation} (generación aumentada por recuperación). \\
\textbf{RAGAS} & \textit{Retrieval Augmented Generation Assessment}, marco de métricas para evaluación de sistemas RAG. \\
\textbf{RAM} & \textit{Random Access Memory} (memoria principal del sistema). \\
\textbf{SLM} & \textit{Small Language Model} (modelo de lenguaje compacto; en este estudio, candidato de hasta 3B parámetros). \\
\textbf{TPS} & \textit{Tokens Per Second} (rendimiento de generación de tokens por segundo). \\
\textbf{VRAM} & \textit{Video Random Access Memory} (memoria gráfica de video dedicada). \\
\bottomrule
\end{tabularx}
\end{table}

% 6. LISTA DE SÍMBOLOS Y NOTACIÓN MATEMÁTICA
\newpage
\phantomsection
\addcontentsline{toc}{chapter}{Lista de símbolos}
\chapter*{Lista de símbolos}

\begin{table}[H]
\centering
\small
\begin{tabularx}{\textwidth}{@{}p{3.2cm}X@{}}
\toprule
\rowcolor{tcAzul!12}
\textbf{Símbolo / Notación} & \textbf{Definición y Rol Computacional} \\
\midrule
$W_0$ & Matriz de pesos preentrenada original del SLM congelada en 4 bits ($W_0 \in \mathbb{R}^{d \times k}$). \\
$\Delta W$ & Matriz de actualización adaptativa acumulada sobre los pesos del modelo ($\Delta W \in \mathbb{R}^{d \times k}$). \\
$A$ & Matriz de proyección de bajo rango inicializada de forma gaussiana ($A \in \mathbb{R}^{r \times k}$). \\
$B$ & Matriz de proyección de bajo rango inicializada en ceros ($B \in \mathbb{R}^{d \times r}$). \\
$r$ & Rango de las matrices de bajo rango, que se fijará durante el diseño experimental ($r \ll \min(d,k)$). \\
$\alpha$ & Factor de escala de la actualización LoRA, cuyo valor se especificará en la configuración experimental. \\
$x$ & Vector de activación de entrada a la capa Transformer (tokens de consulta + contexto RAG). \\
$h$ & Vector de activación adaptado de salida: $h = W_0 x + \frac{\alpha}{r} (B \cdot A) x$. \\
$Q, K, V$ & Matrices de Consulta (\textit{Query}), Clave (\textit{Key}) y Valor (\textit{Value}) en las capas de autoatención. \\
$W_q, W_k, W_v, W_o$ & Matrices de proyección lineal de los cabezales de atención objetivo del ajuste fino QLoRA. \\
$d_k$ & Dimensión de escalamiento de los vectores de clave en el producto escalar de autoatención. \\
$\mathbf{e}_q$ & Vector de incrustación denso (\textit{embedding}) de la consulta formulada por el usuario. \\
$\mathbf{e}_d$ & Vector de incrustación denso del fragmento documental indexado en la base de conocimiento. \\
$q$ & Consulta de entrada al sistema de recuperación. \\
$d$ & Fragmento documental candidato para una consulta. \\
$\text{sim}(q, d)$ & Similitud del coseno normalizada entre los vectores de consulta y fragmento documental. \\
$\mathcal{D}$ & Colección curada de fragmentos documentales bilingües indexados en la base vectorial local. \\
$\mathcal{C}_k$ & Conjunto de los $k$ fragmentos documentales de mayor afinidad semántica recuperados. \\
$P@k$ & Proporción de fragmentos pertinentes en el conjunto de los primeros $k$ resultados recuperados. \\
\bottomrule
\end{tabularx}
\end{table}

% 7. ÍNDICE DE ECUACIONES
\newpage
\phantomsection
\addcontentsline{toc}{chapter}{Índice de ecuaciones}
\chapter*{Índice de ecuaciones}

\begin{table}[H]
\centering
\small
\begin{tabularx}{\textwidth}{@{}p{2.2cm}p{9.0cm}c@{}}
\toprule
\rowcolor{tcAzul!12}
\textbf{Ecuación} & \textbf{Descripción Matemática y Teórica} & \textbf{Página} \\
\midrule
Ecuación (1) & Métrica de similitud del coseno para recuperación densa de fragmentos documentales & \pageref{eq:similitud_coseno} \\
Ecuación (2) & Precisión entre los primeros $k$ fragmentos recuperados (P@k) & \pageref{eq:precision_at_k} \\
Ecuación (3) & Descomposición de bajo rango para adaptación eficiente de pesos (LoRA / QLoRA) & \pageref{eq:descomposicion_lora} \\
Ecuación (4) & Proyecciones lineales de consulta ($Q$), clave ($K$) y valor ($V$) en capas Transformer & \pageref{eq:proyecciones_atencion} \\
Ecuación (5) & Mecanismo de autoatención con producto escalar escalado (\textit{Scaled Dot-Product Attention}) & \pageref{eq:autoatencion} \\
Ecuación (6) & Proyección lineal de salida multicabezal del bloque de atención Transformer & \pageref{eq:salida_atencion} \\
\bottomrule
\end{tabularx}
\end{table}

% 8. ÍNDICE DE ANEXOS
\newpage
\phantomsection
\addcontentsline{toc}{chapter}{Índice de anexos}
\chapter*{Índice de anexos}

\begin{table}[H]
\centering
\small
\begin{tabularx}{\textwidth}{@{}p{2.6cm}p{8.5cm}c@{}}
\toprule
\rowcolor{tcAzul!12}
\textbf{Identificador} & \textbf{Título y Contenido del Anexo} & \textbf{Página} \\
\midrule
Anexo 1 & Matriz de consistencia lógica de la investigación & \pageref{anexo:matriz_consistencia} \\
Anexo 2 & Ficha técnica para registro de corridas experimentales & \pageref{anexo:ficha_tecnica} \\
Anexo 3 & Rúbrica de evaluación de pertinencia pedagógica & \pageref{anexo:rubrica_pedagogica} \\
Anexo 4 & Banco de consultas patrón de evaluación & \pageref{anexo:banco_consultas} \\
\bottomrule
\end{tabularx}
\end{table}

% 9. GLOSARIO DE TÉRMINOS
\newpage
\phantomsection
\addcontentsline{toc}{chapter}{Glosario de términos}
\chapter*{Glosario de términos}

\begin{description}[leftmargin=1.8em, itemsep=0.6em]
\item[\textbf{Aglutinación:}] Proceso morfológico en el que una palabra puede incorporar varios sufijos, cada uno con una función gramatical o discursiva.
\item[\textbf{Alucinación:}] Respuesta generada que contiene afirmaciones no sustentadas por la información disponible o por las fuentes consultadas.
\item[\textbf{Cuantización:}] Representación de los pesos de un modelo con menor precisión numérica para reducir el espacio requerido en memoria; su efecto sobre calidad y velocidad debe medirse.
\item[\textbf{Design Science Research (DSR):}] Enfoque de investigación que diseña, construye y evalúa un artefacto para atender un problema delimitado.
\item[\textbf{Embedding:}] Representación vectorial de una consulta o fragmento documental usada para calcular similitud.
\item[\textbf{Fidelidad al contexto:}] Grado en que las afirmaciones de una respuesta se sustentan en el contexto recuperado, según una métrica o protocolo definido.
\item[\textbf{Generación aumentada por recuperación (RAG):}] Arquitectura que recupera documentos externos y los incorpora como contexto para una respuesta generada.
\item[\textbf{Inferencia local:}] Ejecución del modelo en el equipo del usuario sin enviar la consulta a un servicio remoto durante la generación.
\item[\textbf{NormalFloat de 4 bits (NF4):}] Formato de cuantización de 4 bits descrito como parte del método QLoRA.
\item[\textbf{QLoRA:}] Método que combina un modelo base cuantizado y adaptadores entrenables de bajo rango; su viabilidad depende del modelo, la configuración y el hardware.
\item[\textbf{Quechua Collao:}] Variedad quechua delimitada como objeto de este estudio; los datos deberán conservar esta etiqueta y su procedencia.
\item[\textbf{Small Language Model (SLM):}] Modelo de lenguaje compacto; para esta propuesta se considera como candidato uno de hasta tres mil millones de parámetros.
\end{description}
